\section{Full List of Evaluated Models}
\label{app:models}

Table~\ref{tab:models} lists every model evaluated in this benchmark, with its family, depth, parameter count, and HuggingFace repository, organized into the four groups of Section~\ref{sec:design-models}. The main comparison uses released model instances rather than intermediate training checkpoints. The four groups contain 59, 24, 11, and 24 models, respectively, for Standard LLMs, Memory-Constrained, BabyLM, and Constrained-BabyLM. All 118 models enter every analysis. Architectural depth metadata can include an embedding state in some implementations; the analysis consistently uses the final representation available in the corresponding evaluation.

The Memory-Constrained and Constrained-BabyLM groups (identical small architectures, web-scale vs.\ BabyLM-scale training data) cover several mechanisms: architectures that bias attention toward nearby tokens through a fixed distance penalty, ALiBi \citep{press2021train,devarda2024locally,clark2025linear}; architectures that replace softmax attention with a content-sensitive forget gate \citep{lin2025forgetting} or a stick-breaking process \citep{tan2025scaling}; and hard local $n$-gram windows that can only attend to the two, three, or five preceding tokens regardless of the input \citep{bu2026staticlinearattentionconstraints}. The web-scale group additionally contains CBR-RNN and its variant CBR-RNN-M, a small recurrent architecture with a bounded context representation \citep{timkey2023limited}, and larger web-scale models built on the same mechanisms (BLOOM with ALiBi, the Forgetting Transformer, and MoLM). The small architectures and their BabyLM-scale counterparts were obtained from the public HuggingFace checkpoints released alongside \citet{bu2026staticlinearattentionconstraints}. The BabyLM group spans submissions and baselines from the BabyLM shared tasks, including the organizers' baselines \citep{charpentier2025findings}, a masked-language-modeling variant \citep{edman2025mask}, an instruction-tuned small model \citep{capone2025classit}, a mixture-of-experts architecture \citep{tapaninaho2025moep}, and a model trained with interactive storytelling feedback \citep{martins2025upon}, alongside the earlier BabyLM challenges \citep{warstadt2023cfp,warstadt2023findings,hu2024secondbabylm}, all obtained from their public HuggingFace repositories. The Standard LLM group spans off-the-shelf families with no human-inspired constraint on either architecture or training data: GPT-2 \citep{radford2019gpt2}, GPT-Neo \citep{black2021gptneo}, OPT \citep{zhang2022opt}, Mamba and Mamba-2 \citep{gu2024mamba}, RWKV-4, Pythia \citep{biderman2023pythia}, OLMo-2 \citep{olmo20252olmo2furious}, Qwen3 \citep{qwen3}, Llama-3, Ministral, and the attention baselines of the Forgetting Transformer \citep{lin2025forgetting}, from 14 million to 66 billion parameters and including both Transformer and state-space architectures.

{\setlength{\tabcolsep}{3pt}
\begin{longtable}{@{}p{1.38cm}p{4.9cm}ccc>{\raggedright\arraybackslash}p{2.3cm}@{}}
\caption{All 118 models evaluated in this benchmark, grouped by family within the model groups of Section~\ref{sec:design-models} (rules separate each fine-grained family). Memory-Constrained models trained on BabyLM-scale data (marked by \texttt{babylm} in the repository name) form the Constrained-BabyLM group; all other models listed under Memory-Constrained are trained on web-scale text. The two ``FoX-release baseline'' rows are the matched baselines released with FoX, with ordinary attention and no forget gate (the Pro variant keeps the output gate and key/value shift of FoX-Pro); they carry no memory constraint and are therefore counted as Standard LLMs. Parameter and layer counts follow the originating release (or the loaded config for the memory-constrained and BabyLM checkpoints; see \texttt{neuroai/models/registry.py}). Training-token/word counts are reported only where documented with confidence by the originating source; a blank cell means this was not reliably available rather than zero.} \label{tab:models} \\
\toprule
\textbf{Family} & \textbf{HuggingFace model ID} & \textbf{\#Params} & \textbf{\#Layers} & \shortstack{\textbf{Tr.}\\\textbf{tokens}} & \textbf{Ref.} \\
\midrule
\endfirsthead
\multicolumn{6}{c}{\tablename\ \thetable{} -- continued} \\
\toprule
\textbf{Family} & \textbf{HuggingFace model ID} & \textbf{\#Params} & \textbf{\#Layers} & \shortstack{\textbf{Tr.}\\\textbf{tokens}} & \textbf{Ref.} \\
\midrule
\endhead
\midrule \multicolumn{6}{r}{\textit{continued on next page}} \\
\endfoot
\bottomrule
\endlastfoot
\multicolumn{6}{l}{\textbf{Standard LLMs}} \\
\midrule
\scriptsize GPT-2 & \textcolor{blue}{\ttfamily\fontsize{7.0}{8.4}\selectfont\seqsplit{openai-community/gpt2}} & \scriptsize 124M & \scriptsize 13 & \scriptsize  & \scriptsize \citep{radford2019gpt2} \\
\scriptsize GPT-2 & \textcolor{blue}{\ttfamily\fontsize{7.0}{8.4}\selectfont\seqsplit{openai-community/gpt2-medium}} & \scriptsize 355M & \scriptsize 25 & \scriptsize  & \scriptsize \citep{radford2019gpt2} \\
\scriptsize GPT-2 & \textcolor{blue}{\ttfamily\fontsize{7.0}{8.4}\selectfont\seqsplit{openai-community/gpt2-large}} & \scriptsize 774M & \scriptsize 37 & \scriptsize  & \scriptsize \citep{radford2019gpt2} \\
\scriptsize GPT-2 & \textcolor{blue}{\ttfamily\fontsize{7.0}{8.4}\selectfont\seqsplit{openai-community/gpt2-xl}} & \scriptsize 1.6B & \scriptsize 49 & \scriptsize  & \scriptsize \citep{radford2019gpt2} \\
\midrule
\scriptsize GPT-Neo/J/NeoX & \textcolor{blue}{\ttfamily\fontsize{7.0}{8.4}\selectfont\seqsplit{EleutherAI/gpt-neo-125m}} & \scriptsize 125M & \scriptsize 13 & \scriptsize  & \scriptsize \citep{black2021gptneo} \\
\scriptsize GPT-Neo/J/NeoX & \textcolor{blue}{\ttfamily\fontsize{7.0}{8.4}\selectfont\seqsplit{EleutherAI/gpt-neo-1.3B}} & \scriptsize 1.3B & \scriptsize 25 & \scriptsize  & \scriptsize \citep{black2021gptneo} \\
\scriptsize GPT-Neo/J/NeoX & \textcolor{blue}{\ttfamily\fontsize{7.0}{8.4}\selectfont\seqsplit{EleutherAI/gpt-neo-2.7B}} & \scriptsize 2.7B & \scriptsize 33 & \scriptsize  & \scriptsize \citep{black2021gptneo} \\
\scriptsize GPT-Neo/J/NeoX & \textcolor{blue}{\ttfamily\fontsize{7.0}{8.4}\selectfont\seqsplit{EleutherAI/gpt-j-6b}} & \scriptsize 6B & \scriptsize 29 & \scriptsize 402B & \scriptsize \citep{wang2021gptj} \\
\scriptsize GPT-Neo/J/NeoX & \textcolor{blue}{\ttfamily\fontsize{7.0}{8.4}\selectfont\seqsplit{EleutherAI/gpt-neox-20b}} & \scriptsize 20B & \scriptsize 45 & \scriptsize 472B & \scriptsize \citep{black2022gptneox} \\
\midrule
\scriptsize OPT & \textcolor{blue}{\ttfamily\fontsize{7.0}{8.4}\selectfont\seqsplit{facebook/opt-125m}} & \scriptsize 125M & \scriptsize 13 & \scriptsize 180B & \scriptsize \citep{zhang2022opt} \\
\scriptsize OPT & \textcolor{blue}{\ttfamily\fontsize{7.0}{8.4}\selectfont\seqsplit{facebook/opt-350m}} & \scriptsize 350M & \scriptsize 25 & \scriptsize 180B & \scriptsize \citep{zhang2022opt} \\
\scriptsize OPT & \textcolor{blue}{\ttfamily\fontsize{7.0}{8.4}\selectfont\seqsplit{facebook/opt-1.3b}} & \scriptsize 1.3B & \scriptsize 25 & \scriptsize 180B & \scriptsize \citep{zhang2022opt} \\
\scriptsize OPT & \textcolor{blue}{\ttfamily\fontsize{7.0}{8.4}\selectfont\seqsplit{facebook/opt-2.7b}} & \scriptsize 2.7B & \scriptsize 33 & \scriptsize 180B & \scriptsize \citep{zhang2022opt} \\
\scriptsize OPT & \textcolor{blue}{\ttfamily\fontsize{7.0}{8.4}\selectfont\seqsplit{facebook/opt-6.7b}} & \scriptsize 6.7B & \scriptsize 33 & \scriptsize 180B & \scriptsize \citep{zhang2022opt} \\
\scriptsize OPT & \textcolor{blue}{\ttfamily\fontsize{7.0}{8.4}\selectfont\seqsplit{facebook/opt-13b}} & \scriptsize 13B & \scriptsize 41 & \scriptsize 180B & \scriptsize \citep{zhang2022opt} \\
\scriptsize OPT & \textcolor{blue}{\ttfamily\fontsize{7.0}{8.4}\selectfont\seqsplit{facebook/opt-30b}} & \scriptsize 30B & \scriptsize 49 & \scriptsize 180B & \scriptsize \citep{zhang2022opt} \\
\scriptsize OPT & \textcolor{blue}{\ttfamily\fontsize{7.0}{8.4}\selectfont\seqsplit{facebook/opt-66b}} & \scriptsize 66B & \scriptsize 65 & \scriptsize 180B & \scriptsize \citep{zhang2022opt} \\
\midrule
\scriptsize Mamba & \textcolor{blue}{\ttfamily\fontsize{7.0}{8.4}\selectfont\seqsplit{state-spaces/mamba-130m-hf}} & \scriptsize 130M & \scriptsize 25 & \scriptsize 300B & \scriptsize \citep{gu2024mamba} \\
\scriptsize Mamba & \textcolor{blue}{\ttfamily\fontsize{7.0}{8.4}\selectfont\seqsplit{state-spaces/mamba-370m-hf}} & \scriptsize 370M & \scriptsize 49 & \scriptsize 300B & \scriptsize \citep{gu2024mamba} \\
\scriptsize Mamba & \textcolor{blue}{\ttfamily\fontsize{7.0}{8.4}\selectfont\seqsplit{state-spaces/mamba-790m-hf}} & \scriptsize 790M & \scriptsize 49 & \scriptsize 300B & \scriptsize \citep{gu2024mamba} \\
\scriptsize Mamba & \textcolor{blue}{\ttfamily\fontsize{7.0}{8.4}\selectfont\seqsplit{state-spaces/mamba-1.4b-hf}} & \scriptsize 1.4B & \scriptsize 49 & \scriptsize 300B & \scriptsize \citep{gu2024mamba} \\
\scriptsize Mamba & \textcolor{blue}{\ttfamily\fontsize{7.0}{8.4}\selectfont\seqsplit{state-spaces/mamba-2.8b-hf}} & \scriptsize 2.8B & \scriptsize 61 & \scriptsize 300B & \scriptsize \citep{gu2024mamba} \\
\midrule
\scriptsize Mamba-2 & \textcolor{blue}{\ttfamily\fontsize{7.0}{8.4}\selectfont\seqsplit{AntonV/mamba2-130m-hf}} & \scriptsize 130M & \scriptsize 25 & \scriptsize 300B & \scriptsize \citep{daogu2024transformers} \\
\scriptsize Mamba-2 & \textcolor{blue}{\ttfamily\fontsize{7.0}{8.4}\selectfont\seqsplit{AntonV/mamba2-370m-hf}} & \scriptsize 370M & \scriptsize 49 & \scriptsize 300B & \scriptsize \citep{daogu2024transformers} \\
\scriptsize Mamba-2 & \textcolor{blue}{\ttfamily\fontsize{7.0}{8.4}\selectfont\seqsplit{AntonV/mamba2-780m-hf}} & \scriptsize 780M & \scriptsize 49 & \scriptsize 300B & \scriptsize \citep{daogu2024transformers} \\
\scriptsize Mamba-2 & \textcolor{blue}{\ttfamily\fontsize{7.0}{8.4}\selectfont\seqsplit{AntonV/mamba2-1.3b-hf}} & \scriptsize 1.3B & \scriptsize 49 & \scriptsize 300B & \scriptsize \citep{daogu2024transformers} \\
\scriptsize Mamba-2 & \textcolor{blue}{\ttfamily\fontsize{7.0}{8.4}\selectfont\seqsplit{AntonV/mamba2-2.7b-hf}} & \scriptsize 2.7B & \scriptsize 61 & \scriptsize 300B & \scriptsize \citep{daogu2024transformers} \\
\midrule
\scriptsize Pythia & \textcolor{blue}{\ttfamily\fontsize{7.0}{8.4}\selectfont\seqsplit{EleutherAI/pythia-14m-deduped}} & \scriptsize 14M & \scriptsize 6 & \scriptsize 300B & \scriptsize \citep{biderman2023pythia} \\
\scriptsize Pythia & \textcolor{blue}{\ttfamily\fontsize{7.0}{8.4}\selectfont\seqsplit{EleutherAI/pythia-31m-deduped}} & \scriptsize 31M & \scriptsize 6 & \scriptsize 300B & \scriptsize \citep{biderman2023pythia} \\
\scriptsize Pythia & \textcolor{blue}{\ttfamily\fontsize{7.0}{8.4}\selectfont\seqsplit{EleutherAI/pythia-70m-deduped}} & \scriptsize 70M & \scriptsize 6 & \scriptsize 300B & \scriptsize \citep{biderman2023pythia} \\
\scriptsize Pythia & \textcolor{blue}{\ttfamily\fontsize{7.0}{8.4}\selectfont\seqsplit{EleutherAI/pythia-160m-deduped}} & \scriptsize 160M & \scriptsize 12 & \scriptsize 300B & \scriptsize \citep{biderman2023pythia} \\
\scriptsize Pythia & \textcolor{blue}{\ttfamily\fontsize{7.0}{8.4}\selectfont\seqsplit{EleutherAI/pythia-410m-deduped}} & \scriptsize 410M & \scriptsize 24 & \scriptsize 300B & \scriptsize \citep{biderman2023pythia} \\
\scriptsize Pythia & \textcolor{blue}{\ttfamily\fontsize{7.0}{8.4}\selectfont\seqsplit{EleutherAI/pythia-1b-deduped}} & \scriptsize 1B & \scriptsize 17 & \scriptsize 300B & \scriptsize \citep{biderman2023pythia} \\
\scriptsize Pythia & \textcolor{blue}{\ttfamily\fontsize{7.0}{8.4}\selectfont\seqsplit{EleutherAI/pythia-1.4b-deduped}} & \scriptsize 1.4B & \scriptsize 25 & \scriptsize 300B & \scriptsize \citep{biderman2023pythia} \\
\scriptsize Pythia & \textcolor{blue}{\ttfamily\fontsize{7.0}{8.4}\selectfont\seqsplit{EleutherAI/pythia-2.8b-deduped}} & \scriptsize 2.8B & \scriptsize 33 & \scriptsize 300B & \scriptsize \citep{biderman2023pythia} \\
\scriptsize Pythia & \textcolor{blue}{\ttfamily\fontsize{7.0}{8.4}\selectfont\seqsplit{EleutherAI/pythia-6.9b-deduped}} & \scriptsize 6.9B & \scriptsize 33 & \scriptsize 300B & \scriptsize \citep{biderman2023pythia} \\
\scriptsize Pythia & \textcolor{blue}{\ttfamily\fontsize{7.0}{8.4}\selectfont\seqsplit{EleutherAI/pythia-12b-deduped}} & \scriptsize 12B & \scriptsize 37 & \scriptsize 300B & \scriptsize \citep{biderman2023pythia} \\
\midrule
\scriptsize OLMo-2 & \textcolor{blue}{\ttfamily\fontsize{7.0}{8.4}\selectfont\seqsplit{allenai/OLMo-2-0425-1B}} & \scriptsize 1B & \scriptsize 17 & \scriptsize 4.0T & \scriptsize \citep{olmo20252olmo2furious} \\
\scriptsize OLMo-2 & \textcolor{blue}{\ttfamily\fontsize{7.0}{8.4}\selectfont\seqsplit{allenai/OLMo-2-1124-7B}} & \scriptsize 7B & \scriptsize 32 & \scriptsize 3.9T & \scriptsize \citep{olmo20252olmo2furious} \\
\scriptsize OLMo-2 & \textcolor{blue}{\ttfamily\fontsize{7.0}{8.4}\selectfont\seqsplit{allenai/OLMo-2-1124-13B}} & \scriptsize 13B & \scriptsize 40 & \scriptsize 5.0T & \scriptsize \citep{olmo20252olmo2furious} \\
\midrule
\scriptsize Qwen3 & \textcolor{blue}{\ttfamily\fontsize{7.0}{8.4}\selectfont\seqsplit{Qwen/Qwen3-0.6B-Base}} & \scriptsize 600M & \scriptsize 29 & \scriptsize 36T & \scriptsize \citep{qwen3} \\
\scriptsize Qwen3 & \textcolor{blue}{\ttfamily\fontsize{7.0}{8.4}\selectfont\seqsplit{Qwen/Qwen3-1.7B-Base}} & \scriptsize 1.7B & \scriptsize 29 & \scriptsize 36T & \scriptsize \citep{qwen3} \\
\scriptsize Qwen3 & \textcolor{blue}{\ttfamily\fontsize{7.0}{8.4}\selectfont\seqsplit{Qwen/Qwen3-4B-Base}} & \scriptsize 4B & \scriptsize 37 & \scriptsize 36T & \scriptsize \citep{qwen3} \\
\scriptsize Qwen3 & \textcolor{blue}{\ttfamily\fontsize{7.0}{8.4}\selectfont\seqsplit{Qwen/Qwen3-8B-Base}} & \scriptsize 8B & \scriptsize 37 & \scriptsize 36T & \scriptsize \citep{qwen3} \\
\scriptsize Qwen3 & \textcolor{blue}{\ttfamily\fontsize{7.0}{8.4}\selectfont\seqsplit{Qwen/Qwen3-14B-Base}} & \scriptsize 14B & \scriptsize 41 & \scriptsize 36T & \scriptsize \citep{qwen3} \\
\midrule
\scriptsize Llama-3 & \textcolor{blue}{\ttfamily\fontsize{7.0}{8.4}\selectfont\seqsplit{meta-llama/Llama-3.2-1B}} & \scriptsize 1B & \scriptsize 17 & \scriptsize 9T & \scriptsize  \\
\scriptsize Llama-3 & \textcolor{blue}{\ttfamily\fontsize{7.0}{8.4}\selectfont\seqsplit{meta-llama/Llama-3.2-3B}} & \scriptsize 3B & \scriptsize 29 & \scriptsize 9T & \scriptsize  \\
\scriptsize Llama-3 & \textcolor{blue}{\ttfamily\fontsize{7.0}{8.4}\selectfont\seqsplit{meta-llama/Llama-3.1-8B}} & \scriptsize 8B & \scriptsize 33 & \scriptsize 15T & \scriptsize  \\
\midrule
\scriptsize Ministral & \textcolor{blue}{\ttfamily\fontsize{7.0}{8.4}\selectfont\seqsplit{mistralai/Ministral-3-3B-Base-2512}} & \scriptsize 3B & \scriptsize 29 & \scriptsize  & \scriptsize  \\
\scriptsize Ministral & \textcolor{blue}{\ttfamily\fontsize{7.0}{8.4}\selectfont\seqsplit{mistralai/Ministral-3-8B-Base-2512}} & \scriptsize 8B & \scriptsize 37 & \scriptsize  & \scriptsize  \\
\scriptsize Ministral & \textcolor{blue}{\ttfamily\fontsize{7.0}{8.4}\selectfont\seqsplit{mistralai/Ministral-3-14B-Base-2512}} & \scriptsize 14B & \scriptsize 41 & \scriptsize  & \scriptsize  \\
\midrule
\scriptsize RWKV-4 & \textcolor{blue}{\ttfamily\fontsize{7.0}{8.4}\selectfont\seqsplit{RWKV/rwkv-4-169m-pile}} & \scriptsize 169M & \scriptsize 13 & \scriptsize  & \scriptsize \citep{peng2023rwkv} \\
\scriptsize RWKV-4 & \textcolor{blue}{\ttfamily\fontsize{7.0}{8.4}\selectfont\seqsplit{RWKV/rwkv-4-430m-pile}} & \scriptsize 430M & \scriptsize 25 & \scriptsize  & \scriptsize \citep{peng2023rwkv} \\
\scriptsize RWKV-4 & \textcolor{blue}{\ttfamily\fontsize{7.0}{8.4}\selectfont\seqsplit{RWKV/rwkv-4-1b5-pile}} & \scriptsize 1.5B & \scriptsize 25 & \scriptsize  & \scriptsize \citep{peng2023rwkv} \\
\scriptsize RWKV-4 & \textcolor{blue}{\ttfamily\fontsize{7.0}{8.4}\selectfont\seqsplit{RWKV/rwkv-4-3b-pile}} & \scriptsize 3B & \scriptsize 33 & \scriptsize  & \scriptsize \citep{peng2023rwkv} \\
\scriptsize RWKV-4 & \textcolor{blue}{\ttfamily\fontsize{7.0}{8.4}\selectfont\seqsplit{RWKV/rwkv-4-7b-pile}} & \scriptsize 7B & \scriptsize 33 & \scriptsize  & \scriptsize \citep{peng2023rwkv} \\
\scriptsize RWKV-4 & \textcolor{blue}{\ttfamily\fontsize{7.0}{8.4}\selectfont\seqsplit{RWKV/rwkv-4-14b-pile}} & \scriptsize 14B & \scriptsize 41 & \scriptsize  & \scriptsize \citep{peng2023rwkv} \\
\midrule
\scriptsize FoX-release baseline & \textcolor{blue}{\ttfamily\fontsize{7.0}{8.4}\selectfont\seqsplit{zhixuan-lin/transformer-llama-760m-longcrawl64-48b}} & \scriptsize 760M & \scriptsize 25 & \scriptsize 48B & \scriptsize \citep{lin2025forgetting} \\
\scriptsize FoX-release baseline & \textcolor{blue}{\ttfamily\fontsize{7.0}{8.4}\selectfont\seqsplit{zhixuan-lin/transformer-pro-760m-longcrawl64-48b}} & \scriptsize 760M & \scriptsize 25 & \scriptsize 48B & \scriptsize \citep{lin2025forgetting} \\
\midrule
\multicolumn{6}{l}{\textbf{Memory-Constrained}} \\
\midrule
\scriptsize CBR-RNN & \textcolor{blue}{\ttfamily\fontsize{7.0}{8.4}\selectfont\seqsplit{cclark/cbr-rnn}} & \scriptsize 17.4M & \scriptsize 1 & \scriptsize  & \scriptsize \citep{timkey2023limited} \\
\scriptsize CBR-RNN & \textcolor{blue}{\ttfamily\fontsize{7.0}{8.4}\selectfont\seqsplit{cclark/cbr-rnn-m}} & \scriptsize 17.4M & \scriptsize 1 & \scriptsize  & \scriptsize \citep{timkey2023limited} \\
\midrule
\scriptsize Stick-Breaking & \textcolor{blue}{\ttfamily\fontsize{7.0}{8.4}\selectfont\seqsplit{Lanni-ni/stickbreaking\_pile\_2layer}} & \scriptsize 27.8M & \scriptsize 2 & \scriptsize  & \scriptsize \citep{tan2025scaling} \\
\scriptsize Stick-Breaking & \textcolor{blue}{\ttfamily\fontsize{7.0}{8.4}\selectfont\seqsplit{Lanni-ni/stickbreaking\_pile\_4layer}} & \scriptsize 48.1M & \scriptsize 4 & \scriptsize  & \scriptsize \citep{tan2025scaling} \\
\scriptsize Stick-Breaking & \textcolor{blue}{\ttfamily\fontsize{7.0}{8.4}\selectfont\seqsplit{Lanni-ni/stickbreaking\_babylm\_10m\_2layer}} & \scriptsize 27.8M & \scriptsize 2 & \scriptsize $\sim$10M w. & \scriptsize \citep{tan2025scaling} \\
\scriptsize Stick-Breaking & \textcolor{blue}{\ttfamily\fontsize{7.0}{8.4}\selectfont\seqsplit{Lanni-ni/stickbreaking\_babylm\_10m\_4layer}} & \scriptsize 48.1M & \scriptsize 4 & \scriptsize $\sim$10M w. & \scriptsize \citep{tan2025scaling} \\
\scriptsize Stick-Breaking & \textcolor{blue}{\ttfamily\fontsize{7.0}{8.4}\selectfont\seqsplit{Lanni-ni/stickbreaking\_babylm\_100m\_2layer}} & \scriptsize 27.8M & \scriptsize 2 & \scriptsize $\sim$100M w. & \scriptsize \citep{tan2025scaling} \\
\scriptsize Stick-Breaking & \textcolor{blue}{\ttfamily\fontsize{7.0}{8.4}\selectfont\seqsplit{Lanni-ni/stickbreaking\_babylm\_100m\_4layer}} & \scriptsize 48.1M & \scriptsize 4 & \scriptsize $\sim$100M w. & \scriptsize \citep{tan2025scaling} \\
\midrule
\scriptsize ALiBi & \textcolor{blue}{\ttfamily\fontsize{7.0}{8.4}\selectfont\seqsplit{Lanni-ni/alibi\_pile\_2layer}} & \scriptsize 27.4M & \scriptsize 2 & \scriptsize  & \scriptsize \citep{press2021train} \\
\scriptsize ALiBi & \textcolor{blue}{\ttfamily\fontsize{7.0}{8.4}\selectfont\seqsplit{Lanni-ni/alibi\_pile\_4layer}} & \scriptsize 45.7M & \scriptsize 4 & \scriptsize  & \scriptsize \citep{press2021train} \\
\scriptsize ALiBi & \textcolor{blue}{\ttfamily\fontsize{7.0}{8.4}\selectfont\seqsplit{Lanni-ni/alibi\_babylm\_10m\_2layer}} & \scriptsize 27.4M & \scriptsize 2 & \scriptsize $\sim$10M w. & \scriptsize \citep{press2021train} \\
\scriptsize ALiBi & \textcolor{blue}{\ttfamily\fontsize{7.0}{8.4}\selectfont\seqsplit{Lanni-ni/alibi\_babylm\_10m\_4layer}} & \scriptsize 45.7M & \scriptsize 4 & \scriptsize $\sim$10M w. & \scriptsize \citep{press2021train} \\
\scriptsize ALiBi & \textcolor{blue}{\ttfamily\fontsize{7.0}{8.4}\selectfont\seqsplit{Lanni-ni/alibi\_babylm\_100m\_2layer}} & \scriptsize 27.4M & \scriptsize 2 & \scriptsize $\sim$100M w. & \scriptsize \citep{press2021train} \\
\scriptsize ALiBi & \textcolor{blue}{\ttfamily\fontsize{7.0}{8.4}\selectfont\seqsplit{Lanni-ni/alibi\_babylm\_100m\_4layer}} & \scriptsize 45.7M & \scriptsize 4 & \scriptsize $\sim$100M w. & \scriptsize \citep{press2021train} \\
\midrule
\scriptsize Forgetting-Transf. & \textcolor{blue}{\ttfamily\fontsize{7.0}{8.4}\selectfont\seqsplit{Lanni-ni/forgetting\_pile\_2layer}} & \scriptsize 27.4M & \scriptsize 2 & \scriptsize  & \scriptsize \citep{lin2025forgetting} \\
\scriptsize Forgetting-Transf. & \textcolor{blue}{\ttfamily\fontsize{7.0}{8.4}\selectfont\seqsplit{Lanni-ni/forgetting\_pile\_4layer}} & \scriptsize 45.7M & \scriptsize 4 & \scriptsize  & \scriptsize \citep{lin2025forgetting} \\
\scriptsize Forgetting-Transf. & \textcolor{blue}{\ttfamily\fontsize{7.0}{8.4}\selectfont\seqsplit{Lanni-ni/forgetting\_babylm\_10m\_2layer}} & \scriptsize 27.4M & \scriptsize 2 & \scriptsize $\sim$10M w. & \scriptsize \citep{lin2025forgetting} \\
\scriptsize Forgetting-Transf. & \textcolor{blue}{\ttfamily\fontsize{7.0}{8.4}\selectfont\seqsplit{Lanni-ni/forgetting\_babylm\_10m\_4layer}} & \scriptsize 45.7M & \scriptsize 4 & \scriptsize $\sim$10M w. & \scriptsize \citep{lin2025forgetting} \\
\scriptsize Forgetting-Transf. & \textcolor{blue}{\ttfamily\fontsize{7.0}{8.4}\selectfont\seqsplit{Lanni-ni/forgetting\_babylm\_100m\_2layer}} & \scriptsize 27.4M & \scriptsize 2 & \scriptsize $\sim$100M w. & \scriptsize \citep{lin2025forgetting} \\
\scriptsize Forgetting-Transf. & \textcolor{blue}{\ttfamily\fontsize{7.0}{8.4}\selectfont\seqsplit{Lanni-ni/forgetting\_babylm\_100m\_4layer}} & \scriptsize 45.7M & \scriptsize 4 & \scriptsize $\sim$100M w. & \scriptsize \citep{lin2025forgetting} \\
\midrule
\scriptsize Hard 2-gram & \textcolor{blue}{\ttfamily\fontsize{7.0}{8.4}\selectfont\seqsplit{Lanni-ni/hard\_2gram\_pile\_2layer}} & \scriptsize 15M & \scriptsize 2 & \scriptsize  & \scriptsize \citep{bu2026staticlinearattentionconstraints} \\
\scriptsize Hard 2-gram & \textcolor{blue}{\ttfamily\fontsize{7.0}{8.4}\selectfont\seqsplit{Lanni-ni/hard\_2gram\_pile\_4layer}} & \scriptsize 28.8M & \scriptsize 4 & \scriptsize  & \scriptsize \citep{bu2026staticlinearattentionconstraints} \\
\scriptsize Hard 2-gram & \textcolor{blue}{\ttfamily\fontsize{7.0}{8.4}\selectfont\seqsplit{Lanni-ni/hard\_2gram\_babylm\_10m\_2layer}} & \scriptsize 15M & \scriptsize 2 & \scriptsize $\sim$10M w. & \scriptsize \citep{bu2026staticlinearattentionconstraints} \\
\scriptsize Hard 2-gram & \textcolor{blue}{\ttfamily\fontsize{7.0}{8.4}\selectfont\seqsplit{Lanni-ni/hard\_2gram\_babylm\_10m\_4layer}} & \scriptsize 28.8M & \scriptsize 4 & \scriptsize $\sim$10M w. & \scriptsize \citep{bu2026staticlinearattentionconstraints} \\
\scriptsize Hard 2-gram & \textcolor{blue}{\ttfamily\fontsize{7.0}{8.4}\selectfont\seqsplit{Lanni-ni/hard\_2gram\_babylm\_100m\_2layer}} & \scriptsize 15M & \scriptsize 2 & \scriptsize $\sim$100M w. & \scriptsize \citep{bu2026staticlinearattentionconstraints} \\
\scriptsize Hard 2-gram & \textcolor{blue}{\ttfamily\fontsize{7.0}{8.4}\selectfont\seqsplit{Lanni-ni/hard\_2gram\_babylm\_100m\_4layer}} & \scriptsize 28.8M & \scriptsize 4 & \scriptsize $\sim$100M w. & \scriptsize \citep{bu2026staticlinearattentionconstraints} \\
\midrule
\scriptsize Hard 3-gram & \textcolor{blue}{\ttfamily\fontsize{7.0}{8.4}\selectfont\seqsplit{Lanni-ni/hard\_3gram\_pile\_2layer}} & \scriptsize 15M & \scriptsize 2 & \scriptsize  & \scriptsize \citep{bu2026staticlinearattentionconstraints} \\
\scriptsize Hard 3-gram & \textcolor{blue}{\ttfamily\fontsize{7.0}{8.4}\selectfont\seqsplit{Lanni-ni/hard\_3gram\_pile\_4layer}} & \scriptsize 28.8M & \scriptsize 4 & \scriptsize  & \scriptsize \citep{bu2026staticlinearattentionconstraints} \\
\scriptsize Hard 3-gram & \textcolor{blue}{\ttfamily\fontsize{7.0}{8.4}\selectfont\seqsplit{Lanni-ni/hard\_3gram\_babylm\_10m\_2layer}} & \scriptsize 15M & \scriptsize 2 & \scriptsize $\sim$10M w. & \scriptsize \citep{bu2026staticlinearattentionconstraints} \\
\scriptsize Hard 3-gram & \textcolor{blue}{\ttfamily\fontsize{7.0}{8.4}\selectfont\seqsplit{Lanni-ni/hard\_3gram\_babylm\_10m\_4layer}} & \scriptsize 28.8M & \scriptsize 4 & \scriptsize $\sim$10M w. & \scriptsize \citep{bu2026staticlinearattentionconstraints} \\
\scriptsize Hard 3-gram & \textcolor{blue}{\ttfamily\fontsize{7.0}{8.4}\selectfont\seqsplit{Lanni-ni/hard\_3gram\_babylm\_100m\_2layer}} & \scriptsize 15M & \scriptsize 2 & \scriptsize $\sim$100M w. & \scriptsize \citep{bu2026staticlinearattentionconstraints} \\
\scriptsize Hard 3-gram & \textcolor{blue}{\ttfamily\fontsize{7.0}{8.4}\selectfont\seqsplit{Lanni-ni/hard\_3gram\_babylm\_100m\_4layer}} & \scriptsize 28.8M & \scriptsize 4 & \scriptsize $\sim$100M w. & \scriptsize \citep{bu2026staticlinearattentionconstraints} \\
\midrule
\scriptsize Hard 5-gram & \textcolor{blue}{\ttfamily\fontsize{7.0}{8.4}\selectfont\seqsplit{Lanni-ni/hard\_5gram\_pile\_2layer}} & \scriptsize 15M & \scriptsize 2 & \scriptsize  & \scriptsize \citep{bu2026staticlinearattentionconstraints} \\
\scriptsize Hard 5-gram & \textcolor{blue}{\ttfamily\fontsize{7.0}{8.4}\selectfont\seqsplit{Lanni-ni/hard\_5gram\_pile\_4layer}} & \scriptsize 28.8M & \scriptsize 4 & \scriptsize  & \scriptsize \citep{bu2026staticlinearattentionconstraints} \\
\scriptsize Hard 5-gram & \textcolor{blue}{\ttfamily\fontsize{7.0}{8.4}\selectfont\seqsplit{Lanni-ni/hard\_5gram\_babylm\_10m\_2layer}} & \scriptsize 15M & \scriptsize 2 & \scriptsize $\sim$10M w. & \scriptsize \citep{bu2026staticlinearattentionconstraints} \\
\scriptsize Hard 5-gram & \textcolor{blue}{\ttfamily\fontsize{7.0}{8.4}\selectfont\seqsplit{Lanni-ni/hard\_5gram\_babylm\_10m\_4layer}} & \scriptsize 28.8M & \scriptsize 4 & \scriptsize $\sim$10M w. & \scriptsize \citep{bu2026staticlinearattentionconstraints} \\
\scriptsize Hard 5-gram & \textcolor{blue}{\ttfamily\fontsize{7.0}{8.4}\selectfont\seqsplit{Lanni-ni/hard\_5gram\_babylm\_100m\_2layer}} & \scriptsize 15M & \scriptsize 2 & \scriptsize $\sim$100M w. & \scriptsize \citep{bu2026staticlinearattentionconstraints} \\
\scriptsize Hard 5-gram & \textcolor{blue}{\ttfamily\fontsize{7.0}{8.4}\selectfont\seqsplit{Lanni-ni/hard\_5gram\_babylm\_100m\_4layer}} & \scriptsize 28.8M & \scriptsize 4 & \scriptsize $\sim$100M w. & \scriptsize \citep{bu2026staticlinearattentionconstraints} \\
\midrule
\scriptsize BLOOM & \textcolor{blue}{\ttfamily\fontsize{7.0}{8.4}\selectfont\seqsplit{bigscience/bloom-560m}} & \scriptsize 560M & \scriptsize 25 & \scriptsize  & \scriptsize \citep{workshop2022bloom} \\
\scriptsize BLOOM & \textcolor{blue}{\ttfamily\fontsize{7.0}{8.4}\selectfont\seqsplit{bigscience/bloom-1b1}} & \scriptsize 1.1B & \scriptsize 25 & \scriptsize  & \scriptsize \citep{workshop2022bloom} \\
\scriptsize BLOOM & \textcolor{blue}{\ttfamily\fontsize{7.0}{8.4}\selectfont\seqsplit{bigscience/bloom-1b7}} & \scriptsize 1.7B & \scriptsize 25 & \scriptsize  & \scriptsize \citep{workshop2022bloom} \\
\scriptsize BLOOM & \textcolor{blue}{\ttfamily\fontsize{7.0}{8.4}\selectfont\seqsplit{bigscience/bloom-3b}} & \scriptsize 3B & \scriptsize 31 & \scriptsize  & \scriptsize \citep{workshop2022bloom} \\
\scriptsize BLOOM & \textcolor{blue}{\ttfamily\fontsize{7.0}{8.4}\selectfont\seqsplit{bigscience/bloom-7b1}} & \scriptsize 7.1B & \scriptsize 31 & \scriptsize  & \scriptsize \citep{workshop2022bloom} \\
\midrule
\scriptsize FoX (Forgetting Transf.) & \textcolor{blue}{\ttfamily\fontsize{7.0}{8.4}\selectfont\seqsplit{zhixuan-lin/fox-llama-760m-longcrawl64-48b}} & \scriptsize 760M & \scriptsize 25 & \scriptsize 48B & \scriptsize \citep{lin2025forgetting} \\
\scriptsize FoX (Forgetting Transf.) & \textcolor{blue}{\ttfamily\fontsize{7.0}{8.4}\selectfont\seqsplit{zhixuan-lin/fox-pro-760m-longcrawl64-48b}} & \scriptsize 760M & \scriptsize 25 & \scriptsize 48B & \scriptsize \citep{lin2025forgetting} \\
\midrule
\scriptsize MoLM & \textcolor{blue}{\ttfamily\fontsize{7.0}{8.4}\selectfont\seqsplit{ibm/MoLM-350M-4B}} & \scriptsize 350M & \scriptsize 25 & \scriptsize  & \scriptsize \citep{shen2023moduleformer} \\
\scriptsize MoLM & \textcolor{blue}{\ttfamily\fontsize{7.0}{8.4}\selectfont\seqsplit{ibm/MoLM-700M-4B}} & \scriptsize 700M & \scriptsize 25 & \scriptsize  & \scriptsize \citep{shen2023moduleformer} \\
\scriptsize MoLM & \textcolor{blue}{\ttfamily\fontsize{7.0}{8.4}\selectfont\seqsplit{ibm/MoLM-700M-8B}} & \scriptsize 700M & \scriptsize 49 & \scriptsize  & \scriptsize \citep{shen2023moduleformer} \\
\midrule
\multicolumn{6}{l}{\textbf{BabyLMs}} \\
\midrule
\scriptsize BabyLM & \textcolor{blue}{\ttfamily\fontsize{7.0}{8.4}\selectfont\seqsplit{colinglab/CLASS\_IT-140M}} & \scriptsize 140M & \scriptsize 13 & \scriptsize  & \scriptsize \citep{capone2025classit} \\
\scriptsize BabyLM & \textcolor{blue}{\ttfamily\fontsize{7.0}{8.4}\selectfont\seqsplit{llm-slice/blm-gpt2s-90M-s42\_901M-s42\_submission}} & \scriptsize 90M & \scriptsize 13 & \scriptsize  & \scriptsize  \\
\scriptsize BabyLM & \textcolor{blue}{\ttfamily\fontsize{7.0}{8.4}\selectfont\seqsplit{Jtapsa/moep}} & \scriptsize 80M & \scriptsize 13 & \scriptsize  & \scriptsize \citep{tapaninaho2025moep} \\
\scriptsize BabyLM & \textcolor{blue}{\ttfamily\fontsize{7.0}{8.4}\selectfont\seqsplit{BabyLM-community/BabyLM-2026-Baseline-GPT2-Strict}} & \scriptsize 98.4M & \scriptsize 13 & \scriptsize $\sim$100M w. & \scriptsize \citep{charpentier2025findings} \\
\scriptsize BabyLM & \textcolor{blue}{\ttfamily\fontsize{7.0}{8.4}\selectfont\seqsplit{BabyLM-community/BabyLM-2026-Baseline-GPT2-Strict-Small}} & \scriptsize 98.4M & \scriptsize 13 & \scriptsize $\sim$10M w. & \scriptsize \citep{charpentier2025findings} \\
\scriptsize BabyLM & \textcolor{blue}{\ttfamily\fontsize{7.0}{8.4}\selectfont\seqsplit{BabyLM-community/babylm-baseline-100m-gpt-bert-causal-focus}} & \scriptsize 119.5M & \scriptsize 13 & \scriptsize $\sim$100M w. & \scriptsize \citep{edman2025mask} \\
\scriptsize BabyLM & \textcolor{blue}{\ttfamily\fontsize{7.0}{8.4}\selectfont\seqsplit{BabyLM-community/babylm-baseline-100m-gpt-bert-mixed}} & \scriptsize 119.5M & \scriptsize 13 & \scriptsize $\sim$100M w. & \scriptsize \citep{edman2025mask} \\
\scriptsize BabyLM & \textcolor{blue}{\ttfamily\fontsize{7.0}{8.4}\selectfont\seqsplit{BabyLM-community/babylm-baseline-100m-gpt-bert-masked-focus}} & \scriptsize 119.5M & \scriptsize 13 & \scriptsize $\sim$100M w. & \scriptsize \citep{edman2025mask} \\
\scriptsize BabyLM & \textcolor{blue}{\ttfamily\fontsize{7.0}{8.4}\selectfont\seqsplit{BabyLM-community/babylm-baseline-10m-gpt-bert-causal-focus}} & \scriptsize 33M & \scriptsize 13 & \scriptsize $\sim$10M w. & \scriptsize \citep{edman2025mask} \\
\scriptsize BabyLM & \textcolor{blue}{\ttfamily\fontsize{7.0}{8.4}\selectfont\seqsplit{BabyLM-community/babylm-baseline-10m-gpt-bert-mixed}} & \scriptsize 33M & \scriptsize 13 & \scriptsize $\sim$10M w. & \scriptsize \citep{edman2025mask} \\
\scriptsize BabyLM & \textcolor{blue}{\ttfamily\fontsize{7.0}{8.4}\selectfont\seqsplit{BabyLM-community/babylm-baseline-10m-gpt-bert-masked-focus}} & \scriptsize 33M & \scriptsize 13 & \scriptsize $\sim$10M w. & \scriptsize \citep{edman2025mask} \\
\end{longtable}
}

\section{Full List of Evaluated Datasets}
\label{app:datasets}

Table~\ref{tab:datasets} lists every dataset, grouped by modality, corresponding to the behavioral and neural categories of Section~\ref{sec:design-data}. The 15 datasets yield 16 behavioral measures and eight neural datasets; behavioral measures from one corpus are separate outcomes, not independent replications.

{\setlength{\tabcolsep}{3pt}
\begin{longtable}{@{}p{2.45cm}p{1.5cm}cp{2.6cm}p{1.8cm}>{\raggedright\arraybackslash}p{2.0cm}@{}}
\caption{All datasets evaluated in this benchmark, grouped by modality. $N$ items reports the unit each dataset's items are counted in (sentences, passages, screens, stories, electrodes, \ldots), since this varies by design. Eye-tracking measures: FP = first-pass/gaze duration, GP = go-past time, SP = total dwell time. $N$ subj.\ reflects each study's usable sample after its own quality-control exclusions.}
\label{tab:datasets} \\
\toprule
\textbf{Dataset} & \textbf{Modality} & \textbf{Subj.} & \textbf{Items} & \textbf{Measure(s)} & \textbf{Ref.} \\
\midrule
\endfirsthead
\multicolumn{6}{c}{\tablename\ \thetable{} -- continued} \\
\toprule
\textbf{Dataset} & \textbf{Modality} & \textbf{Subj.} & \textbf{Items} & \textbf{Measure(s)} & \textbf{Ref.} \\
\midrule
\endhead
\midrule \multicolumn{6}{r}{\textit{continued on next page}} \\
\endfoot
\bottomrule
\endlastfoot
\multicolumn{6}{l}{\textbf{Self-Paced Reading}} \\
\midrule
\small Natural Stories SPR & \small SPR & \small 180 & \small 10 stories (10{,}256 word pos.) & \small raw RT & \scriptsize \citep{futrell2021natural} \\
\small UCL Reading & \small SPR & \small 117 & \small 361 sentences & \small raw RT & \scriptsize \citep{frank2013reading} \\
\midrule
\multicolumn{6}{l}{\textbf{Eye-Tracking}} \\
\midrule
\small GECO & \small Eye-tracking & \small 14 & \small 588 screens & \small FP, GP & \scriptsize \citep{cop2017geco} \\
\small MECO (English) & \small Eye-tracking & \small 46 & \small 12 passages & \small SP, FP, GP & \scriptsize \citep{siegelman2022meco} \\
\small Provo & \small Eye-tracking & \small 84 & \small 55 texts & \small SP, FP, GP & \scriptsize \citep{luke2018provo} \\
\small OneStop & \small Eye-tracking & \small 360 & \small 330 passages & \small SP, FP, GP & \scriptsize \citep{berzak2025onestop} \\
\small ZuCo Reading & \small Eye-tracking & \small 12 & \small SR/NR/TSR: 400/300/407 sent. & \small TRT & \scriptsize \citep{hollenstein2018zuco} \\
\midrule
\multicolumn{6}{l}{\textbf{N400 / ERP}} \\
\midrule
\small Nieuwland & \small EEG (ERP) & \small 334 & \small 9-lab RSVP replication & \small N400 & \scriptsize \citep{nieuwland2018largescale} \\
\small Federmeier et al.\ (2007) & \small EEG (ERP) & \small 32 & \small 564 items & \small N400 & \scriptsize \citep{federmeier2007multiple,szewczyk2022context} \\
\small Hubbard et al.\ (2019) & \small EEG (ERP) & \small 32 & \small 192 items & \small N400 & \scriptsize \citep{szewczyk2022context} \\
\midrule
\multicolumn{6}{l}{\textbf{Continuous EEG}} \\
\midrule
\small Broderick & \small EEG & \small 19 & \small 20 segs. ($\sim$196s each) & \small EEG (1--8Hz) & \scriptsize \citep{broderick2018electrophysiological} \\
\midrule
\multicolumn{6}{l}{\textbf{MEG}} \\
\midrule
\small Gwilliams (MEG-MASC) & \small MEG & \small 27 & \small 4 stories $\times$ 2 sess. & \small continuous MEG & \scriptsize \citep{gwilliams2022neural} \\
\midrule
\multicolumn{6}{l}{\textbf{fMRI}} \\
\midrule
\small Tuckute & \small fMRI & \small 10 & \small 2{,}000 sentences (6 ROIs) & \small BOLD & \scriptsize \citep{tuckute2024driving} \\
\small Pereira & \small fMRI & \small $\leq$10 & \small 627 sentences (12 ROIs) & \small BOLD & \scriptsize \citep{pereira2018toward} \\
\midrule
\multicolumn{6}{l}{\textbf{iEEG / ECoG}} \\
\midrule
\small Podcast & \small ECoG & \small 9 & \small 1 story, $>$5{,}000 w. (1{,}330 electr.) & \small high-gamma & \scriptsize \citep{zada2025podcast} \\
\end{longtable}
}

\section{Implementation Details}
\label{app:implementation}

\paragraph{Data provenance.} All datasets were obtained from their public OpenNeuro or OSF repositories under the terms specified therein.

\paragraph{Feature extraction.} All language models are loaded in bfloat16 and evaluated as frozen feature extractors, with no gradient updates. For each word, the hidden states of every layer are mean-pooled over its constituent subword tokens, and a word's representation is computed with its full preceding sentence or discourse context in the forward pass rather than in isolation. Word surprisal is the sum of the subword-token surprisals of the model's own output distribution. Preprocessing, splits and confounds are identical across models.

\paragraph{Matched-scale regressions.} The all-datasets RSA regressions (Equation~\ref{eq:group}; Tables~\ref{tab:summary}, \ref{tab:summary-z}, \ref{tab:rsa-sameq} and \ref{tab:rsa-ckpt}) are fitted with \texttt{lme4} by restricted maximum likelihood (BOBYQA optimizer), with dataset-unit and lineage random intercepts and Satterthwaite $t$-tests (\texttt{lmerTest}). For the main fits (Table~\ref{tab:summary}, all datasets), the estimated standard deviations are $.209$ (dataset unit), $.021$ (lineage) and $.079$ (residual) for behavioral, and $1.683$, $.072$ and $.788$ for neural similarity; with robust-$z$ outcomes they are $.170$, $.292$, $.821$ and $.260$, $.383$, $1.313$. None of these fits is singular. The word-pair bin regressions (Table~\ref{tab:summary}) and the network regressions (Appendix~\ref{app:md}) use the same estimator and the lineage random intercept, with lineage-by-bin and, for the 9 reading measures, measure-by-bin intercepts in the bin fits and dataset fixed effects in the network fits. Earlier versions fitted the bins and networks by ordinary least squares with family- or model-clustered standard errors; the coefficients are close, but the clustered standard errors were larger for the DLT bins (memory constraint on any-hard pairs $+.017$, $p=.13$, against $+.018$, $p=.001$ here) and smaller for the networks. The two-dataset analysis (Appendix~\ref{app:two-datasets}) is an ordinary least-squares fit with standard errors clustered by model family; the per-measure frequency fits (Table~\ref{tab:freq-measures}) are ordinary least-squares fits with HC3 standard errors.



